\chapter{Introduzione}
\label{ch:introduzione}

Questo elaborato valuta empiricamente se e quando un DBMS distribuito sia adatto a un carico applicativo reale di tipo HR/retributivo (ditte, dipendenti, contratti, presenze e cedolini paga in regime
\emph{multitenant}), confrontando due sistemi che rappresentano filosofie opposte di gestione dei dati:

\begin{itemize}
  \item \textbf{PostgreSQL + Citus}: un DBMS \emph{relazionale} reso realmente distribuito dall'estensione
        Citus (un coordinator e $N$ worker), con schema rigido, SQL completo, transazioni ACID e consistenza
        forte;
  \item \textbf{MongoDB}: un DBMS \emph{documentale} (sharded cluster), con schema flessibile, dati
        organizzati in documenti annidati e distribuzione automatica tramite il balancer.
\end{itemize}

I due sistemi appartengono alla stessa classe rispetto al teorema CAP (entrambi \textbf{CP}) ma hanno modelli di dati molto diversi: questo permette di capire quando ciascun modello rappresenta bene questo tipo di dati e come si comportano in ambiente distribuito.

\section{Contesto e motivazione}
\label{sec:usecase}

La scelta del dominio non è casuale: nasce da una collaborazione lavorativa con un'azienda di consulenza HR,
tra i cui servizi principali c'è la creazione e l'elaborazione dei cedolini paga. Partendo da \textbf{file XML
di paghe reali}, è sorta una domanda concreta: come modellare al meglio questi dati, sia con un database
\emph{relazionale} sia con uno \emph{documentale} (non relazionale). In un contesto realistico con
\textbf{migliaia di dipendenti e di aziende}, si tratta poi di capire quale dei due sistemi si comporti meglio
e quindi quale convenga implementare in produzione.

Prima di avviare i test sono stati studiati i fondamenti teorici dei due sistemi e lette le rispettive
documentazioni ufficiali: il modello dei dati, il linguaggio di interrogazione e il modo in cui affrontano
scalabilità e transazioni concorrenti. Solo dopo si è passati alla verifica sperimentale. L'elaborato segue
questa struttura: una parte teorica (\Cref{ch:dominio}) e una sperimentale su un cluster reale
(\Cref{ch:metodologia,ch:risultati}), in cui si confronta quanto affermato in teoria con i dati misurati. Per il vincolo GDPR i dati reali non sono utilizzabili direttamente: il dataset è
quindi sintetico ma \emph{calibrato} sulla forma dei cedolini reali (\Cref{sec:dati}).

\section{Requisiti e criteri di valutazione}
\label{sec:requisiti}

Il lavoro è organizzato attorno ai requisiti richiesti. Occorre disporre di un dataset ampio (o reso
artificialmente grande) da interrogare con query a complessità variabile su entrambi i sistemi, e valutare
\emph{quando il modello modella bene} il dominio (\Cref{ch:progettazione}). Vanno poi verificati la gestione
dei cambiamenti di schema in corso d'opera (\Cref{sec:evoluzione}), il comportamento in ambiente distribuito
in lettura, scrittura (transazioni) e in presenza di guasti, e la scalabilità nei due regimi: a parità di
dati aggiungendo nodi, e a parità di nodi aumentando il volume dei dati. Questi ultimi aspetti sono misurati
sperimentalmente (\Cref{ch:metodologia,ch:risultati}).

\section{Approccio e struttura}
\label{sec:approccio}

Il metodo è sperimentale: entrambi i sistemi sono stati installati su un cluster di 5 macchine virtuali in
cloud, alimentati con lo \emph{stesso dataset} sintetico (calibrato su cedolini reali) e sollecitati con una
suite di query identica nella semantica. Per ogni prova si sono raccolte latenza, throughput e uso di
risorse (CPU, memoria, I/O di rete e disco) \emph{di ciascun nodo}, così da osservare non solo \emph{quanto}
costa un'operazione, ma \emph{come} il lavoro si distribuisce nel cluster.
